\documentclass[11pt]{article}

\usepackage[margin=1in]{geometry}
\usepackage{amsmath}
\usepackage{amssymb}
\usepackage{mathtools}
\usepackage{enumitem}
\usepackage{parskip}
\usepackage{graphicx}
\usepackage{listings}
\usepackage{xcolor}
\usepackage{hyperref}

\lstdefinestyle{pythonstyle}{
  language=Python,
  basicstyle=\ttfamily\small,
  keywordstyle=\color{blue!70!black}\bfseries,
  stringstyle=\color{green!50!black},
  commentstyle=\color{gray}\itshape,
  numberstyle=\tiny\color{gray},
  numbers=left,
  numbersep=8pt,
  showstringspaces=false,
  breaklines=true,
  frame=single,
  framesep=4pt,
  xleftmargin=14pt,
  tabsize=4,
}

\newcommand{\flops}{\text{FLOPs}}
\newcommand{\flopss}{\text{FLOPs/s}}
\newcommand{\bytesps}{\text{bytes/s}}
\newcommand{\AI}{\mathrm{AI}}
\newcommand{\Tmath}{T_{\text{math}}}
\newcommand{\Tcomm}{T_{\text{comm}}}

\title{Scaling Book --- Section 1 Exercises}
\author{Rahul Bir}
\date{}

\begin{document}
\maketitle

\section*{Q1}

Let $X \in \mathbb{R}^{B,D}$, $Y \in \mathbb{R}^{D,F}$, $Z \in \mathbb{R}^{B,F}$, with
\[
XY = Z \quad \text{in int8.}
\]

\subsection*{(i) Bytes loaded / written}
\begin{align*}
\text{loaded:} &\quad BD + DF \\
\text{written:} &\quad BF
\end{align*}

\subsection*{(ii) FLOPs}
There are $BF$ dot products, each of cost $D + D - 1$:
\[
\underbrace{BF}_{\text{dot prods}} \cdot \underbrace{(D + D - 1)}_{\text{dot prod cost}} \approx 2BFD.
\]

\subsection*{(iii) Arithmetic intensity}
\[
\AI = \frac{2BFD}{BD + DF + BF}.
\]
If $B \ll D, F$:
\[
\AI \approx \frac{2BFD}{DF} = 2B.
\]

\subsection*{(iv) Lower / upper bound on runtime}
\[
\Tmath = \frac{\flops}{\flopss} = \frac{2BFD}{3.94\mathrm{e}14},
\qquad
\Tcomm = \frac{\text{bytes}}{\bytesps} = \frac{BD + DF + BF}{8.2\mathrm{e}11}.
\]
\begin{align*}
\text{lower bound:} &\quad \max(\Tmath,\, \Tcomm) \\
\text{upper bound:} &\quad \Tmath + \Tcomm
\end{align*}
(plug in above)

\section*{Q2}

\[
\text{bf16}[B,D] \;\cdot\; \text{int8}[D,F] \;\longrightarrow\; \text{bf16}[B,F]
\]

\begin{align*}
\text{total } \flops &= 2BFD \\
\text{bytes moved} &= 2BD + DF + 2BF
\end{align*}

\[
\AI = \frac{2BFD}{2BD + DF + 2BF} \quad (\text{assume } B \ll F, D) \;=\; 2B.
\]

At the corner of the roofline model:
\[
\text{Intensity}(\text{accel}) = \frac{1.97\mathrm{e}14 \; \flops}{8.2\mathrm{e}11 \; \text{bytes}}.
\]

\[
2B = \text{Intensity}(\text{accel})
\quad\Longrightarrow\quad
B = \frac{1.97\mathrm{e}14}{2 \cdot 8.2\mathrm{e}11} \approx 120.
\]

\section*{Q3}

Done on computer. Plot peak achievable FLOPs/s as a function of batch size for $F = D = 4096$ and $F = D = 1024$ (TPU v5e: $1.97 \times 10^{14}$ \flopss, $8.2 \times 10^{11}$ \bytesps).

\lstinputlisting[style=pythonstyle]{q3.py}

\begin{figure}[h]
  \centering
  \includegraphics[width=0.75\linewidth]{q3_graph.png}
  \caption*{Peak FLOPs/s vs.\ batch size. Blue: $F = D = 4096$. Orange: $F = D = 1024$. Both saturate at the accelerator peak ($1.97 \times 10^{14}$ \flopss); the larger matmul ($D = 4096$) reaches the corner of the roofline at a smaller batch size.}
\end{figure}

\section*{Q4}

\[
\text{int8}[B,D] \;\cdot\; \text{int8}[B,D,F] \;\longrightarrow\; \text{int8}[B,F]
\]

(Per-batch matmul: for each $b$, a length-$D$ vector multiplies a $D \times F$ matrix.)

\begin{align*}
\flops &= BF(2D) = 2DBF \\
\text{bytes} &= BD + BDF + BF \\
\AI &= \frac{2DBF}{BD + BDF + BF}
\end{align*}

\section*{Q5}

H100 SXM specs:
\begin{align*}
\text{bf16:} &\quad \underbrace{1.979 \times 10^{15} \; \flopss}_{\text{w/ sparsity}}
\;\Longrightarrow\; \text{true value: } \frac{1.979}{2} \times 10^{15} \; \flopss \\
\text{bw:} &\quad 3.35 \times 10^{12} \; \bytesps
\end{align*}

$\AI(\text{matmul}) = B$ \quad (assuming small $B$).

\[
B_{\text{compute-bound}}
= \left\lfloor \frac{\frac{1.979}{2} \times 10^{15} \; \flopss}{3.35 \times 10^{12} \; \bytesps} \right\rfloor
\approx 295 \text{ tokens} \quad (\text{rounding down}).
\]

\end{document}
